\chapter{相关工作}

本章从多模态大语言模型、并行与分布式训练基础、异构 GPU 集群并行训练策略以及
梯度检查点优化等方面对国内外相关研究进行综述。
如今，基于 Transformer 的大语言模型在包括自然语言处理\cite{brown2020gpt3,openai2023}、
图像识别\cite{vit2020}、多模态\cite{clip2021}、图学习与推荐系统等众多前沿领域和
下游任务上大放异彩。
然而，随着模型参数的增加和数据规模的扩大，模型结构越来越复杂、
模型参数量越来越大、训练数据越来越多，
导致单个设备已经无法完成模型和数据的存储，
多模态大语言模型的训练越来越有难度。
为了解决上述问题，目前学界已经提出并应用了多种并行策略来实现和优化大模型的训练。
然而，目前传统的并行策略往往忽略了并行训练平台的异构性，
由于服务器集群中常常存在多种型号的 GPU 和网络连接，
导致集群中存在着存储异构、算力异构和网络异构。
在这样的复杂集群环境下，
传统的并行策略会因为缺乏对异构环境的感知和适应而导致性能表现急剧下滑。
因此，下文将主要从现有大模型的并行训练策略出发，
对并行训练的研究现状进行详细论述，并讨论各个并行策略在异构环境下的优缺点。

\section{多模态大语言模型}

\subsection{多模态大语言模型架构}

多模态大语言模型（MLLMs）通常由 Modality Encoder、Input Projector、LLM Backbone、
Output Projector 等模块组成。
Modality Encoder 负责将不同模态的输入数据（如图像、视频、音频）编码为模型可理解的表示，
常采用预训练的视觉模型如 CLIP\cite{clip2021}、ViT\cite{vit2020} 等。
LLM Backbone 是模型的核心，通常为参数量巨大的 Transformer 基座模型（如 GPT 系列%
\cite{brown2020gpt3,gpt2}、LLaMA\cite{llama3}、T5\cite{t5} 等），
占据整个模型约 80\%--90\% 的参数量。
Modality Generator 一般为 Diffusion 类生成模型（如 Stable Diffusion\cite{stable_diffusion}），
根据 LLM 的输出生成图像、视频等其它模态数据。
前后两个 Projector 用于将不同模态映射到共享语义空间，
可能是 Linear 层、Transformer 层或 Q-Former\cite{blip2} 等结构。
Blip-2\cite{blip2} 率先提出冻结预训练 Encoder 与 LLM、仅训练轻量 Projector 的范式，
后续 NextGPT\cite{nextgpt}、DreamLLM\cite{dreamllm} 等工作均沿用并发展了此类架构。
多模态大语言模型的这种模块异构性使得不同层在算力、显存和通信上的需求差异很大，
在异构 GPU 集群上部署时面临独特的挑战。

\subsection{多模态大语言模型训练数据}

为获得满意的性能，大语言模型与多模态大语言模型通常依赖大规模数据训练。
例如图像数据集 JFT-3B 包含约 30 亿张图片，
Common Crawl 网页数据包含数千亿 tokens，
Open Image\cite{open_images} 需约 18TB 存储，
YouTube-8M\cite{youtube8m} 等视频数据集甚至需要 PB 级存储。
在单个硬件加速器上既无法存储完整模型参数，也无法容纳如此大规模的训练数据，
因此并行训练（数据并行、模型并行及其组合）成为必然选择。
数据规模与模型规模的持续增长，
也使得对并行策略的通信效率、存储划分和异构适应性提出了更高要求。

\subsection{多模态大语言模型训练目标}

多模态大语言模型的训练通常分为多个阶段，如模态对齐、指令微调等，
不同阶段会冻结不同模块（如冻结 LLM 仅训练 Projector，
或冻结 Encoder 与 LLM 仅训练部分层）。
冻结会显著改变各模块的反向计算量与显存占用（如优化器状态仅需为未冻结参数维护），
并影响流水线中各 stage 的负载均衡。
在异构 GPU 集群下，这种因阶段和冻结策略带来的“模型异构”与设备异构叠加，
进一步增加了并行策略设计与自动调优的难度。

\section{并行与分布式训练基础}

\subsection{常见并行策略}

数据并行（Data Parallelism）被广泛用于处理数据规模过大的情况，
其核心思想是通过将训练数据在所有计算设备上划分，
而每个计算设备都持有一份完整的模型参数\cite{parameter_server,horovod,pytorch_ddp}。
如图\ref{fig:data_parallel}所示，在训练过程的每次迭代中，
每个设备可以独立地使用本地模型参数和批次数据计算对应批次的梯度；
然后，所有设备在下一次迭代之前进行通信操作，
将各自计算得到的梯度数据进行同步，
从而保证模型参数的一致性和结果的正确性。
数据并行训练因其隐式性和直观性而成为并行训练最早提出的策略，
也是实际应用中最常见的方式。
PyTorch 框架中提供了 DistributedDataParallel\cite{pytorch_ddp} 工具用于数据并行训练。
Horovod\cite{horovod} 将数据并行策略适配到了多个框架中。
Parallax\cite{parallax} 则根据模型的特点灵活地使用参数服务器（Parameter Server）或
All-Reduce 通信原语进行数据同步，优化了数据并行中的通信开销。
然而，数据并行存在着明显限制：每个 GPU 都保存模型权重的完整副本，
因此不能直接用于训练具有大量参数的模型，限制了模型规模的扩展；
同时，数据并行也将模型训练的工作限制在了单个计算设备内，
无法通过设备并行的方式加速训练过程。

\begin{figure}[htbp]
  \centering
  \includegraphics[width=0.75\linewidth]{fig03_data_parallel}
  \caption{数据并行}
  \label{fig:data_parallel}
\end{figure}

模型并行（Model Parallelism）是实现大模型训练的有效途径，
包括张量并行（Tensor Parallelism）\cite{megatron2020} 和
流水线并行（Pipeline Parallelism）\cite{gpipe}。
张量并行的核心思想是将模型在层内进行划分，
每个计算设备分别持有模型一部分的层内张量，它们组合在一起是一份完整的张量，
从而将模型的参数和计算任务分别绑定到不同的相互连接的设备上。
如图\ref{fig:tensor_parallel}所示，在训练过程的每次迭代中，
每个计算设备只负责模型一部分参数的计算，
然后通过中间激活值和梯度值的通信，使模型可以进行下一层的计算。
Megatron-LM\cite{megatron2020} 提出了针对 Transformer 模型张量并行的训练方法，
对每两个连续层的权重首先按行（输入维度）划分，再按列（输出维度）划分，
消除了同步第一层中间输出的需要，但随后仍然需要大量的跨设备通信。
GShard\cite{gshard} 提出了在计算图层面进行张量并行优化的思路。
然而，目前的张量并行方式受限于 GPU 服务器集群中的带宽情况，尤其是跨节点带宽，
其并行规模无法进行大规模扩展，往往需要结合其他并行方式一起加速。

\begin{figure}[htbp]
  \centering
  \includegraphics[width=0.75\linewidth]{fig04_tensor_parallel}
  \caption{张量并行}
  \label{fig:tensor_parallel}
\end{figure}

流水线并行的核心思想是将模型在层间进行划分，
每个计算设备分别持有整个模型的若干层，称为一个 stage；
这些 stage 组合起来是一份完整的模型。
如果直接将不同设备上的 stage 串行地进行训练，
则会导致每个 stage 训练过程中只有一个设备在进行计算，设备利用率极低。
如图\ref{fig:pipeline_parallel}所示，流水线并行针对训练数据进一步划分为多个
micro-batch，依次输入到模型中进行计算。
尽管各级流水线在不同的 micro-batch 之间存在数据依赖，
然而 batch 之间数据是相互独立的，
从而每个 micro-batch 可以依次进行模型训练，实现流水线的训练过程，
大大减少计算设备的空闲时间。
GPipe\cite{gpipe} 提出了使用 micro-batch 的流水线并行方式，
通过设计多级流水减少了 Bubble 的空闲。
PipeDream\cite{pipedream2019} 提出了 1F1B（1 次前向 1 次反向），
通过异步调度的方式进一步减少了 Bubble 的空闲。
TeraPipe\cite{terapipe} 则在自回归模型的单个训练 sequence 中使用 token 级流水线并行，
实现了更细粒度的流水线并行。
数据并行、张量并行和流水线并行策略相互正交，可以一起用于并行训练。
由于流水线并行的通信量一般较小，且跨节点通信明显慢于节点内通信，
因此张量并行往往仅限于节点内训练，而流水线并行则用于跨节点训练，
以更好地利用硬件带宽\cite{megatron2020}。

\begin{figure}[htbp]
  \centering
  \includegraphics[width=0.8\linewidth]{fig05_pipeline_parallel}
  \caption{流水线并行}
  \label{fig:pipeline_parallel}
\end{figure}

\subsection{多维混合并行}

ZeRO\cite{zero2020} 是微软 DeepSpeed 团队提出的零冗余优化器技术，
针对大规模深度神经网络中数据并行训练进行高效优化。
如图\ref{fig:zero}所示，在模型训练中，
各部分的存储占用从高到低依次为优化器状态、梯度和参数。
ZeRO 通过将以上三部分分别划分到每个计算设备中，
进而减少存储开销和提高 GPU 利用率来扩展大语言模型的训练。
在训练过程中，ZeRO 通过前向和反向的多次通信同步数据，保证训练结果的一致性。
目前，学界已经提出了一些针对 ZeRO 的通信方面的改进工作。
MiCS\cite{mics2022} 是在云服务环境下针对 ZeRO 的通信数据量进行减少和优化；
ZeRO++\cite{zeroplusplus} 基于量化技术减少 ZeRO 的通信数据量；
AMSP\cite{amsp} 则将 ZeRO 的划分方式扩展到一个统一划分空间，
寻找通信开销最小的策略，进一步扩展了大模型训练的规模。

\begin{figure}[htbp]
  \centering
  \includegraphics[width=0.8\linewidth]{fig06_zero}
  \caption{零冗余优化器 ZeRO}
  \label{fig:zero}
\end{figure}

\section{异构 GPU 集群并行训练策略}

在存储方面，由于不同 GPU 之间的存储不一致，
传统的并行策略会导致存储负载不均衡，
大存储的设备会出现存储利用不充分的问题，影响了模型训练的扩展性。
在算力方面，由于不同 GPU 之间的算力不一致，
传统并行策略会导致“木桶效应”，不仅使得计算资源利用率下降，
还导致整体训练速度的降低。
在网络方面，由于集群中存在的不对称的多级网络连接，
使得不同设备之间的通信开销相去甚远，
低带宽网络对应的通信将成为训练的性能瓶颈。
因此，面向异构 GPU 集群的并行训练策略成为近年来的研究热点。

HetPipe\cite{hetpipe2020} 结合了流水线并行（PMP）和数据并行（DP），
多个 GPU 被分组为虚拟工作节点（VW），以流水线方式处理小批量数据，
然后这些虚拟工作节点再通过数据并行提高训练的可扩展性和性能。
但是 HetPipe 没有考虑虚拟工作节点内部的异构设备排列方式，
对于通信开销和显存开销的考虑过于直观，而且完全没有考虑模型异构，
在如今的多模态大语言模型训练时不适用。

Whale\cite{whale} 是阿里巴巴团队提出的一种面向异构存储的 GPU 集群的并行训练优化方法。
由于实际 GPU 集群中存在的存储异构和算力异构问题，
Whale 提出了一种基于硬件感知的负载均衡算法，
通过按比例分配的方式为每个设备分配不同大小的批次数据和模型部分，
实现了数据并行与张量并行的负载均衡。
然而，Whale 没有考虑集群中普遍存在的网络异构性，
对于存储异构和算力异构的优化方法过于直观和简单，
并且是在子图的粗粒度下进行的划分和编译，也没有考虑 ZeRO 的并行方法，
因此仍然需要进一步的优化。

AMP\cite{amp2022} 是卡耐基梅隆大学团队提出的一种针对异构环境下大模型的
并行训练优化方法。
AMP 同时考虑了模型异构和设备异构，
针对模型中不同层的结构不同的情况进行了抽象和建模，
而针对集群中不同 GPU 存储异构和网络异构进行了近似计算。
在此基础上，AMP 基于 3D 并行方法（DP+TP+PP），
设计了一个详细的开销模型（Cost Model）以评估异构模型在异构环境下的计算和通信开销，
并通过动态规划算法确定最优的并行策略。
尽管 AMP 在传统的 3D 并行策略的基础上提升了异构环境下的性能表现，
但是 AMP 的并行维度是基于全局的粗粒度划分，开销模型的准确性不足，
并且没有考虑显存开销，
在实际的应用中会大概率发生内存不足（Out of Memory）的问题，
故而仍然具有很大的改进空间。

而与 ZeRO 相关的工作中，不论是 ZeRO、MiCS 还是 ZeRO++ 和 AMSP，
在面向大语言模型并行训练的异构性环境下，均存在着潜在的问题和提升的空间。
ZeRO 中存在着庞大的通信开销；
MiCS 中存在着对模型规模的限制，并且对于异构的网络连接中的高带宽连接利用率低；
ZeRO++ 的量化方法引入了潜在的信息损失和模型性能的下降，存在通用性问题；
AMSP 针对通信的优化思路同样局限于减少跨节点的通信量，
无法充分适应和利用异构集群中存在的不同存储、算力和网络带宽。

\section{梯度检查点优化}

在大规模模型与长序列训练中，为节省显存，
常采用梯度检查点（Gradient Checkpointing）技术，
即不保存所有层的中间激活，仅在部分层保存，
反向传播时对未保存的层进行重新前向计算以得到梯度。
全部重计算策略不保存任何中间激活，仅在需要时按需重算，
显存占用最低但计算量增加明显。
选择重计算策略则选择部分层作为检查点保存激活，在显存与计算之间取得折中。
细粒度重计算进一步在算子或更小粒度上决定保存与重算，
可结合模型结构和异构设备的显存与算力进行更精细的优化。
在异构 GPU 集群上，不同设备显存与算力差异大，
将梯度检查点与流水线划分、stage 分配结合，
可避免显存溢出并提高整体吞吐，是本文后续优化方向之一。
Colossal-Auto\cite{colossal_auto} 等工作对自动并行与激活检查点进行了统一自动化，
为面向异构环境的联合优化提供了参考。
